\appendix
\section{Derivation chain and limiting controls}
The record calculation begins with a normalized joint state and a declared controlled unitary. Partial trace gives the off-diagonal overlap. A balanced analyser turns that overlap into visibility. Optimizing a specified marker decoder gives trace-norm distinguishability. Purification and contraction extend the pure-state equality to mixed states and restricted access. The Born rule and the quantum channel formalism are inputs throughout.

The restoration calculation begins with a complete unitary and a specified initial recipient. Inner-product preservation constrains the output records. A three-register permutation realizes one reversible write-and-transfer operation. It makes the local cell return while retaining its correlation in history; it also shows why an occupied recipient cannot be treated as a fresh blank cell. A correlated source can instead uncompute a write without adding a new history register. These are distinct protocols.

The original-clock chain begins with the supplied exact hard eigenbranches and soft vacuum. A matched reference retains the even energy and soft counterterm. The only difference is the square-integrable soft interaction. Its action on the reference preparation lies in the one-soft-photon sector, orthogonal to the reference energy component. The form representation supplies the full operator-domain statement; the squared norm gives the variance correction. The same soft norm, constant along the reference orbit, gives the independent Duhamel overlap enclosure. This chain does not infer a full spectral measure from a single clock current.

Useful limiting controls are: identical versus orthogonal markers; biased versus balanced priors; conditioned versus unsorted eraser data; marker-only reset versus a joint controlled inverse; blank versus occupied history memory; equal versus unequal route proper ages; a common scalar energy reassignment with versus without its compensating external phase; and an omitted soft coupling versus the retained full-Hamiltonian remainder. Each rejected substitution changes a stated operand or experiment.

\section{Reproduction and data availability}
The companion bundle contains this manuscript, its figure generator, the retained Chapter I model and certificate inputs, pinned quantum-route/history source snapshots, and independent mathematical checks. The new numerical work comprises reproducible finite-state tests and exact rational propagation of the inherited six-lane intervals. Floating-point finite comparators corroborate the displayed identities within their recorded tolerances; they do not replace proofs or certify a truncated model as the original continuum.

The original Chapter I and II public repositories remain their own evidence sources.\cite{chapter1,chapter2} The Chapter IV and V inputs used here are the retained expanded manuscripts and supplied source/replay materials. No Chapter VI repository URL has been assigned. The bundled provenance manifest identifies all input bytes and pinned repository paths; the included reproduction command runs without a native GEN6 engine or network access.

The CR075 source package is a historical structural audit. Twenty-eight recovered repository files reproduce their pinned Git blob identities. Historical CR075 checksum differences reconcile as line-ending differences without rewriting the retained Git snapshots. This draft does not report a fresh CR075, QP, or native GEN6 campaign. Inherited interacting spectral certificates are authenticated inputs; their original continuum/native computation is not re-executed by the new finite tests.

The declared hard-eigenbranch preparation is a mathematical state specification, not an achieved laboratory preparation. Similarly, the proper-time controlled unitary, allowed full-marker measurement, ideal recombination and decoder are receiving assumptions. Completing their physical implementation is part of the source-linked work listed above.

